% Generated draft: requires mathematical and source review.
\documentclass[11pt]{article}
\usepackage[margin=25mm]{geometry}
\usepackage{fontspec}
\setmainfont{Latin Modern Roman}
\usepackage{amsmath,amssymb,mathtools,unicode-math}
\setmathfont{Latin Modern Math}
\usepackage{enumitem,needspace,xurl,hyperref}
\setlength{\parindent}{0pt}
\setlength{\parskip}{5pt}
\setlength{\emergencystretch}{3em}
\newcommand{\N}{\mathbb{N}}
\newcommand{\Z}{\mathbb{Z}}
\newcommand{\Q}{\mathbb{Q}}
\newcommand{\R}{\mathbb{R}}
\newcommand{\C}{\mathbb{C}}
\newcommand{\F}{\mathbb{F}}
\newcommand{\A}{\mathbb{A}}
\newcommand{\PP}{\mathbb{P}}
\newcommand{\E}{\mathbb{E}}
\newcommand{\eps}{\varepsilon}
\newcommand{\dd}{\,\mathrm{d}}
\begin{document}
\section*{UnsolvedMath ID: 20000421}
\textbf{Source problem code:} catalogue entry
\subsection{Definitions and mathematical statement}
C. Schoen asks to describe the surfaces in $J\times J$ for a general abelian surface $J$, over $\C$ for example. A surface means a two-dimensional closed subvariety; possible targets include individual embedded surfaces, Chow classes, numerical classes, or the effective/pseudo-effective cone. The original does not choose among them or specify a polarization type and the meaning of ``general.'' Its surrounding numerals ``0,'' ``3'' and ``1'' are unexplained extraction artifacts; no different product or verified source numbering is inferred from them.

The attached unreviewed aid restricts to a very general principally polarized complex abelian surface $(J,\Theta)$. Let $p_i:J\times J\to J$, let $\theta_i=p_i^*[\Theta]$, and fix the Poincar\'e convention $\lambda=m^*[\Theta]-\theta_1-\theta_2$, where $m$ is addition and a symmetric polarization representative is used. In the real codimension-two numerical cycle space it reports the basis
\[
\theta_1^2,\ \theta_1\theta_2,\ \theta_2^2,\
\theta_1\lambda,\ \theta_2\lambda,\ \lambda^2.
\]
For a class $\alpha$, set
\[
(X_1,\ldots,X_6)=
(\alpha\theta_1^2,\alpha\theta_1\theta_2,\alpha\theta_2^2,
 \alpha\theta_1\lambda,\alpha\theta_2\lambda,\alpha\lambda^2),
\]
where products mean intersection degrees. The aid attributes an exact pseudo-effective cone to Debarre--Ein--Lazarsfeld--Voisin and claims recovery of coefficients by inversion of the intersection matrix. Its proposed test requires
\[
6X_2+X_6\ge2|X_2+X_6|
\]
and positivity of a specified scaled $3\times3$ matrix. That matrix and the inversion coefficients are only referred to as equations (5.8) of an unavailable \texttt{report.md}; they are not supplied here. Thus the displayed inequality alone is not presented as a complete criterion. Pseudo-effective means closure of the cone generated by effective numerical surface classes.

In the same very-general setting, the aid claims every abelian subgroup surface is the image of $x\mapsto(mx,nx)$ for a primitive integer pair $(m,n)$, with class
\[
\frac12(n^2\theta_1+m^2\theta_2-mn\lambda)^2.
\]
A primitive pair has $\gcd(m,n)=1$; simultaneous sign change gives the same subgroup. Translating such a subgroup gives further embedded surfaces with its numerical class. These numerical and subgroup descriptions do not classify all embedded surfaces or their Chow classes, and do not extend without qualification from very general principally polarized surfaces to the original unspecified family.
\subsection{Short English statement}
Clarify whether describing surfaces in an abelian-square means geometric classification or cycle-cone computation. The aid gives numerical data and subgroup classes only for a very general principally polarized surface; its complete positivity matrix is missing.
\subsection{Sources}
\begin{itemize}
\item C. Schoen, supplied item ``0,'' AIM \emph{The Tate conjecture}: \url{https://aimath.org/WWN/tateconjecture/tateconjecture.pdf}. Canonical location: \texttt{aim-algebraic-number-theory-notes.json}, notes[72]. The record is labeled exact despite unexplained surrounding numerals; remarks and literature are empty. No PDF was consulted.
\item The attached aid invokes Debarre--Ein--Lazarsfeld--Voisin and equations (5.8) of \texttt{report.md}, without supplying the matrix, full theorem or bibliographic entry. It is machine generated and not peer reviewed, with low novelty confidence.
\item Local registry: \texttt{data/unsolvedmath/problems.json}, permanent catalogue entry, exact source code \texttt{catalogue entry}.
\end{itemize}
\end{document}
